\clearpage

\section{Summary of changes included in version 3 of T2K-TN-154 (released on , 2013)}

This section is designed to allow readers to keep track of the changes between versions of T2K-TN-154.
Table, figure, equation, footnote or citation numbers shown in black refer to version of 3,
{\color{red}whereas the numbers shown in red refer to version of 4}.
Both numbers are shown where possible.

\begin{itemize}

\item{
Abstract changed to be more accurate in the following way:  "where all possible, but currently unknown, correlations" $\rightarrow$ "where possible, but currently unknown, correlations". \ck
}

\item{
Text to explain the check marks was added to the caption in Table 1 ({\color{red}2}). \ck
}

\item{
An extra reference was added, TN34, ({\color{red}[6]}) to give more details about SK efficiencies for $1R_{\mu}$. \ck
}

\item{
\textit{" Indicate the solar oscillation parameter values that were used.
"}
Added in Table {\color{red}1}. \ck
}

\item{
\textit{"Fix the unusual numerical labeling of $sin^{2}\theta_{23}$ for figures on page 10. (Also for figures 36, 37, 39-52).
"}
\ck
}

\item{
\textit{"167: The LHS depends on the 83 systematic parameters. This is not indicated in the notation - which could cause confusion. 
"}
The dependency on the systematic parameters is indicated in the term S which depends on a vector of nuisance (systematic) parameters $\vec{f}$. \ck
}


\item{
\textit{"167: Both r and r' are referred to as indices for SK reconstructed energy bins. The distintion between them is not made until the next paragraph, line 185, and even then, it is not explicit. Perhaps close to equation (4) it could be stated that r' is the reconstructed SK energy bin in absence of a shift in energy scale. Line 185 is incorrect: replace `due to uncertainty in the SuperK reconstructed energy scale� by `due to a potential shift in the SuperK reconstructed energy scal�
"}

Sentence added to explain the difference between r and r'. Line corrected. \ck
}


\item{
\textit{"Equations (4) and (5) have very different notation for the expected number of events. Can this notation be unified, or if the variables are not identical, can an equation relating them be given?
"}

Notation in equation (5) changed to match equation (4). \ck
}


\item{
\textit{"244 ({\color{red}246}): remove ( before "is performed" 
"}\ck
}


\item{
\textit{"248 ({\color{red}250}): add [ before "-5 sigma" 
"} \ck
}


\item{
\textit{"Table 11 and Figures 14,17,18,19 are not referred to in the text.
"}
References have been added in the text. \ck
}


\item{
\textit{"Inconsistent naming: sometimes "single ring", sometimes "1 ring", sometimes "one ring".
"}

Unified to "single ring" except when abbreviations are used (1R$\mu$ and 1R_{e}). \ck
}


\item{
\textit{"356 ({\color{red}358}): !.2 should be 1.2
"}\ck
}

\item{
\textit{"Figure 39 and 40: should indicate that dashed in 68\% and solid is 90\%. Define the middle solid red line in figure 40. 
"}
\ck
}

\item{
\textit{" 536 ({\color{red}538}): replace "tere" by "there"
"}
\ck
}

\end{itemize}






